\section*{Overview}

Binary search is not really about sorted arrays. The broader pattern is: there is an ordered search space, and once a
candidate becomes valid, everything on one side stays valid as well. Then I can replace a linear scan with repeated
halving.

In contests this appears in several forms: searching inside a sorted container, finding the first time a process
becomes feasible, and binary searching an answer when the real work is hidden inside a checker.

\section*{When to Use It}

Use binary search when:

\begin{itemize}
  \item the search space is ordered,
  \item the predicate is monotone,
  \item checking one candidate is much cheaper than trying every candidate,
  \item the final answer is a boundary such as ``first valid'', ``last valid'', or an approximate real value.
\end{itemize}

Typical examples are minimum capacity, maximum achievable score, smallest time, or first position where a prefix
condition becomes true.

\section*{Core Idea}

Maintain an interval that is guaranteed to contain the answer. At each step, evaluate the middle point and throw away
the half that cannot contain the boundary anymore.

For a first-true search on integers, I keep the invariant:
\[
\texttt{answer} \in [lo, hi], \qquad \texttt{pred(hi) = true}.
\]
If \(\texttt{pred(mid)}\) is true, the answer is in the left half. Otherwise it is in the right half.

\section*{Key Insight}

The main job is not writing the loop. The main job is choosing a predicate whose truth set is monotone.

For example, if I want the minimum value \(x\) such that some task can be completed, I do not search directly on the
construction. I ask a yes/no question: ``Can I do it with limit \(x\)?'' If the answer changes from false to true only
once, the problem is ready for binary search.

\section*{Operations / Main Technique}

\paragraph{Sorted array search.}
Use the standard lower-bound / upper-bound interpretation:
\begin{itemize}
  \item first position with value \(\ge x\),
  \item first position with value \(> x\),
  \item last position with value \(\le x\).
\end{itemize}

\paragraph{Binary search on answer.}
Pick a candidate answer \(mid\), run a greedy, DP, or prefix-based checker, and use the result as a monotone
predicate.

\paragraph{Real-valued search.}
If the answer is a real number, run a fixed number of iterations instead of relying on exact equality.

\paragraph{Worked example.}
Suppose I need the smallest machine speed that finishes all jobs within \(T\) hours. If speed \(s\) works, then every
speed larger than \(s\) also works. That turns the optimization problem into a first-true boundary search.

\section*{Correctness Intuition}

Each iteration preserves the invariant that the answer still lies in the active interval. The monotonicity of the
predicate guarantees that once one side is rejected, no valid answer is lost with it.

When the interval shrinks to one point, that point must be the boundary we were maintaining.

\section*{Complexity Analysis}

If the search space has size \(N\), integer binary search needs \(O(\log N)\) predicate calls. The total cost is
\[
O(\log N \cdot \text{cost of checker}).
\]
For real-valued search, the complexity is \(O(k \cdot \text{cost of checker})\), where \(k\) is the chosen iteration
count.

\section*{Implementation}

The reference code includes:

\begin{itemize}
  \item \texttt{first\_true} on integers,
  \item \texttt{last\_true} on integers,
  \item a floating-point variant with a fixed number of iterations.
\end{itemize}

I prefer writing these as small templates and keeping the predicate outside. That makes the checker easy to test on
its own.

\section*{Common Pitfalls}

\begin{itemize}
  \item Using binary search without proving the predicate is monotone.
  \item Picking bounds that do not actually contain the answer.
  \item Mixing ``first true'' and ``last true'' logic in the same loop.
  \item Overflow in \(\texttt{mid = (lo + hi) / 2}\) when the bounds are large.
  \item Using binary search on real values when the problem really wants an exact discrete argument.
\end{itemize}

\section*{Variants / Extensions}

\begin{itemize}
  \item Exponential search to discover an upper bound before binary search.
  \item Parallel binary search for many offline monotone queries.
  \item Binary lifting style search on Fenwick trees or doubling tables.
  \item Ternary search when the target is unimodal rather than monotone.
\end{itemize}

\section*{Practice Problems}

\begin{itemize}
  \item Find the first position of a value in a sorted array with duplicates.
  \item Minimum capacity to ship packages within \(D\) days.
  \item Maximum minimum distance after placing objects greedily.
  \item Binary search on answer where the checker uses prefix sums or greedy feasibility.
\end{itemize}

\section*{References}

\begin{itemize}
  \item \href{https://cp-algorithms.com/num_methods/binary_search.html}{cp-algorithms: Binary Search}.
  \item \href{https://usaco.guide/silver/binary-search}{USACO Guide: Binary Search}.
  \item \href{https://codeforces.com/catalog}{Codeforces Catalog}.
\end{itemize}
